
% !TeX root = ../main.tex
\begin{englishabstract}
    \footnotetext{*This work was supported by Fundamental and Interdisciplinary Disciplines Breakthrough Plan of the Ministry of Education of China (JYB2025XDXM116), National Natural Science Foundation of China (No.62450005,62477036).}
    % \astfootnote{The work was supported by the Foundation (foundation ID).}

    Plane geometry problem solving is a typical multimodal task integrating visual understanding, text parsing, and logical reasoning, which has important application value in intelligent tutoring and educational digitalization. Although current multimodal large models have achieved remarkable breakthroughs in multimodal understanding and basic mathematical reasoning, they still have critical performance limitations in long-chain geometric reasoning tasks. Models generally suffer from the lack of evaluation benchmarks, visual noise interference, and cumulative errors in intermediate reasoning. Especially in scenarios of multi-step continuous deduction, models are prone to reasoning deviation and fail to maintain stable and reliable long-chain deduction capabilities. To address the above issues, this thesis conducts systematic research on long-chain plane geometry reasoning. The main contributions and innovations are as follows:

    To tackle the lack of high-quality benchmarks and incomplete process evaluation for long-chain geometric reasoning, this thesis constructs a multimodal plane geometry benchmark named GeoLStep. The benchmark collects 2097 high-difficulty comprehensive geometry problems from real senior high school entrance examinations, covering numerical calculation and conclusion proof tasks. It provides fine-grained annotations including reasoning steps and auxiliary lines, and builds a full-pipeline image-text formalization conversion framework. This thesis establishes a dual-dimensional evaluation protocol based on answer accuracy and step correctness, and conducts comparative experiments on multiple baseline models. Experimental results show that the performance of mainstream models degrades by up to 54.5\% in 13--26 step long-chain reasoning, which fully verifies the core challenges for existing models and provides standardized evaluation and formalization support for complex geometric reasoning research.

    To solve the core problems of logical instability and error accumulation in long-chain geometric reasoning, this thesis proposes a two-stage enhancement method GeoTSR combining a weak symbolic verifier and group relative policy optimization. The method adopts a modality decoupling training strategy: it first strengthens the basic deduction ability of models in a text-only environment, and then introduces step-level rule feedback in multimodal scenarios to suppress error propagation through validity checking and redundancy elimination. Experimental results demonstrate that the proposed method outperforms traditional Chain-of-Thought and supervised fine-tuning baselines, with an average accuracy improvement of 7.0\% on the MathVerse benchmark, as well as 7.3\% higher answer accuracy and 17.0\% higher step correctness on the long-chain subset of GeoLStep, which significantly improves the stability and accuracy of long-chain reasoning.

    This thesis designs and implements an interactive plane geometry problem-solving prototype system, which integrates dataset management, multi-model parallel scheduling, and reasoning result visualization. The system supports multimodal input, heterogeneous algorithm comparison, and interpretable result output, verifying the engineering feasibility and practical value of the proposed methods in intelligent teaching scenarios.

    \englishkeywordstype{Plane Geometry Problem Solving; Multimodal Reasoning; Long-Chain Reasoning; Benchmark Construction; Reinforcement Learning}{Application Research}

\end{englishabstract}